\documentclass{exam}

\newif\ifA
\newif\ifkey
\newif\ifprintselected   % required by exam_config's \ansbox

%%%%%%%%%%%%%%%%%%%%%%%%%%%%%%%%%%%%%%%%%%%%%%%%%%%%%%
%%% SET THESE IF VARIABLES
%%%%%%%%%%%%%%%%%%%%%%%%%%%%%%%%%%%%%%%%%%%%%%%%%%%%%%
% Lab name -- labs are named by the lesson they belong to, never by a
% sequential lab number (Neil, 2026-10-12): this is the L11 Lab.
\def\labLesson{11}

% is this the key or not?
\ifdefined\iskey
	\ifnum\iskey=1
		\keytrue
	\else
		\keyfalse
	\fi
\else
	%% Pick one manually
	\keyfalse % if this is NOT the key
%	\keytrue     % if this IS the key
\fi

% set up the header and footer - change for every term
\ifdefined\thisTermIn
	\def\thisterm{\thisTermIn}
\else
	\def\thisterm{Spring 2027}
\fi
%
\def\thisclass{ECE 448}

%%%%%%%%%%%%%%%%%%%%%%%%%%%%%%%%%%%%%%%%%%%%%%%%%%%%%%
%%% END SET THESE IF VARIABLES
%%%%%%%%%%%%%%%%%%%%%%%%%%%%%%%%%%%%%%%%%%%%%%%%%%%%%%
% packages and shortcuts
\input{myPackages_gr}
\usepackage[hidelinks, linkcolor=red]{hyperref}
\input{myShortcuts}
\setlength{\parindent}{0pt}
\setlength{\parskip}{6pt}

%% exam class customization
\input{exam_config}
% ECE 448 document fonts.
%
% Body text is Barlow (SIL OFL; the faces and the license are in
% latex/fonts/Barlow/). Barlow is not in TeX Live, so it is loaded as an
% OpenType face through fontspec -- which means the document must be built
% with lualatex, not pdflatex. build_practice.sh and build_lab.sh both do.
%
% Compiling with pdflatex anyway is not an error: the guard below falls back
% to the LaTeX default so the source still compiles under any engine and can
% be syntax-checked. The resulting PDF just is not the one we ship.
%
% Paths are relative to latex/, which is where both build scripts run
% (each does cd "$(dirname "$0")"). Loading by path rather than by installed
% family name means no system font install and no fontconfig, so a fresh
% clone builds a byte-comparable PDF anywhere.
%
% MATH IS DELIBERATELY LEFT IN COMPUTER MODERN. Barlow has no math companion,
% and CM math under sans body text is already the house pairing: the decks set
% words in Source Sans Pro and let MathJax render symbols in its CM-derived
% face. Moving math to a sans as well would mean unicode-math plus a sans math
% font (Fira Math being the plausible one) and re-checking every equation in
% the course for symbol coverage. That is a separate decision, not a side
% effect of this one.

\usepackage{iftex}
\ifLuaTeX
	\usepackage{fontspec}
	\defaultfontfeatures{Ligatures=TeX}
	\setmainfont{Barlow}[
		Path           = fonts/Barlow/,
		Extension      = .ttf,
		UprightFont    = *-Regular,
		ItalicFont     = *-Italic,
		BoldFont       = *-Bold,
		BoldItalicFont = *-BoldItalic,
	]
	\setsansfont{Barlow}[
		Path           = fonts/Barlow/,
		Extension      = .ttf,
		UprightFont    = *-Regular,
		ItalicFont     = *-Italic,
		BoldFont       = *-SemiBold,
		BoldItalicFont = *-BoldItalic,
	]
	% microtype earns its place here specifically because of Barlow. Its
	% glyphs are narrower than Computer Modern's, so justified prose sets
	% tighter and TeX has to squeeze interword space to fit; two paragraphs
	% in the L02/L03 sets went overfull on the first Barlow build. Protrusion
	% and font expansion absorb that -- both are lualatex features, so this
	% belongs inside the guard with everything else.
	\usepackage{microtype}

	% Mono stays Latin Modern: Barlow ships no monospaced face, and verbatim
	% and code need one that actually is fixed-pitch.
\fi
   % Barlow body text (needs lualatex)
\printselectedfalse      % full worked solutions in the key, not answers-only
\CorrectChoiceEmphasis{\color{red}}
\SolutionEmphasis{\color{red}}
\renewcommand\questionlabel{\textbf{\thequestion.}}
\renewcommand{\thepartno}{\alph{partno}}
\renewcommand{\partlabel}{(\thepartno)}

% Fillable table cell: blank in the student copy, the expected value in red on
% the KEY. Fixed-width centered columns keep the blank cells writable.
\newcolumntype{K}[1]{>{\centering\arraybackslash}p{#1}}
\newcommand{\kv}[1]{\ifkey{\color{red}#1}\fi\rule[-0.10in]{0pt}{0.34in}}

% print the answers if this is the key
\ifkey
	\printanswers % if this IS the key
\else
	\noprintanswers  %if this is NOT the key
\fi

% print the header and footer
% need to print key on top if this is the key
\ifkey
	\firstpageheader{}{\color{red}{ KEY}}{}
	\runningheader{\thisclass\hspace{1pt} -- \thisterm }{L\labLesson\ Lab\hspace{1pt}\color{red}{ KEY}}{\thepage\hspace{1pt} of\hspace{1pt} \numpages}
	\firstpagefooter{}{}{}
	\runningfooter{}{}{}
\else
	\firstpageheader{}{}{}
	\runningheader{\thisclass \hspace{1pt} -- \thisterm }{L\labLesson\ Lab\hspace{1pt}}{\thepage\hspace{1pt} of\hspace{1pt} \numpages}
	\firstpagefooter{}{}{}
	\runningfooter{}{}{}
\fi

\runningheadrule


\begin{document}
\pagestyle{headandfoot}   % exam class

%%%%%%%%%%%%%%%%%%%%%%%%%%%%%%%%%%%%%%%%%%%%%%%%%%%%%%
% title block
%%%%%%%%%%%%%%%%%%%%%%%%%%%%%%%%%%%%%%%%%%%%%%%%%%%%%%

\begin{center}
USAF ACADEMY \\
DEPARTMENT OF ELECTRICAL AND COMPUTER ENGINEERING \\
LESSON \labLesson\ LAB - CUSTOM BLOCK \\
\textbf{Lab Packet} \\
\thisterm \\[4mm]
\end{center}

\vspace{-4pt}
\textbf{Name:} \rule[-2pt]{2.6in}{0.4pt} \hfill
\textbf{Section:} \rule[-2pt]{0.9in}{0.4pt} \hfill
\textbf{Date:} \rule[-2pt]{1.1in}{0.4pt}

\vspace{4pt}

\fullwidth{\textbf{Learning Objective 5.1: } I can create a GNU Radio out-of-tree (OOT) module
and add a custom Python block to it.
\\[2pt]\textbf{Learning Objective 5.2: } I can implement a custom signal-processing block from
scratch and integrate it into a GNU Radio flowgraph.}

\vspace{2pt}

This packet is the L\labLesson\ Lab, in three parts, and you turn it in once, complete,
before the midterm. Do \textbf{Part 1} on your own before the lab period: check your
laptop's tools and make the predictions you will check at the bench. \textbf{Part 2} is
your own work in class: build and test the course module, set the Energy Detector's
levels, and then write a block of your own, the \textbf{Key Timer}, in a module of your
own. Do \textbf{Part 3} on your own after the lab.

\textbf{Grading.} One submission, graded against Objectives 5.1 and 5.2 on the course's
1 / 0 scale, with revision allowed under course policy. Paste screenshots and code where
a question asks for them.

Nothing in this lab transmits: every flowgraph is a simulation on your laptop.

The files are on the L11 lesson page: \texttt{L11-gr-ece448.zip},
\texttt{L11-keyed-carrier.grc}, and \texttt{L11-keyed-carrier-epy.grc}. Every number
below assumes the flowgraph's settings: \texttt{samp\_rate} = 32 kHz, a tone of
amplitude 1, Morse ``ECE'' repeating, and GNU Radio's Noise Source, whose power is
$\sigma^2$ at amplitude $\sigma$. One Morse unit lasts $1.2/\text{WPM}$ seconds: a dit
is 1 unit, a dah 3, the gap inside a letter 1, the gap between letters 3, and the gap
after a word 7.

\textbf{Two routes.} On Linux or macOS with \texttt{gnuradio-build-deps} installed, you
build real modules. On a Windows laptop, take the
Embedded Python Block route the lab page describes: put the same code in Embedded Python
Blocks, and build the two modules on a teammate's Linux or macOS laptop for Questions 3(a)
and 5.

\begin{questions}

\fullwidth{\large\textbf{Part 1: Pre-Lab} \normalsize(individual, before the lab period)}

%%%%%%%%%%%%%%%%%%%%%%%%%%%%%%%%%%%%%%%%%%%%%%%%%%%
\question \textbf{Your tools.} At the radioconda prompt, run
\texttt{gnuradio-config-info --version}, \texttt{gr\_modtool --help}, and
\texttt{cmake --version}.

\begin{parts}

	\begin{minipage}[t]{0.58\linewidth}
		\part Record your operating system, your GNU Radio version, and whether
		\texttt{gr\_modtool} and \texttt{cmake} run. Name the route you will take.
		\begin{solution}[0.9in]
		Radioconda installs GNU Radio 3.10 with \texttt{gr\_modtool} on every platform.
		\texttt{cmake} runs only after \texttt{conda install -c conda-forge
		gnuradio-build-deps}. Linux and macOS take the module route; Windows takes the
		Embedded Python Block route.
		\end{solution}
	\end{minipage}\hfill%
	\begin{minipage}[t]{0.37\linewidth}
		\raggedleft \vspace{12pt}
		\ansbox[1.6in]{3.10.x; route}
	\end{minipage}

	\part Open and run \texttt{L11-keyed-carrier-epy.grc}, which needs nothing installed.
	Paste a screenshot of its window at the default settings.
		\begin{solution}[1.5in]
		The power steps between about 0 dB and $-20$ dB in the pattern of ECE, and the
		detected key matches the sent keying.
		\end{solution}

\end{parts}

%%%%%%%%%%%%%%%%%%%%%%%%%%%%%%%%%%%%%%%%%%%%%%%%%%%
\newpage
\question \textbf{Predictions.} You will check each of these at the bench.

\begin{parts}

	\part The bench uses a noise amplitude of 0.5. Predict the Energy Detector's power
	output with the carrier off and on, and a threshold halfway between them in dB.
	Repeat for a noise amplitude of 0.8.

	\begin{center}
	\renewcommand{\arraystretch}{1.25}
	\begin{tabular}{| l | K{1.1in} | K{1.1in} | K{1.1in} |}
	\hline
	$\sigma$ & Carrier off (dB) & Carrier on (dB) & Threshold (dB) \\\hline
	0.5 & \kv{$-6.0$} & \kv{$+1.0$} & \kv{$-2.5$} \\\hline
	0.8 & \kv{$-1.9$} & \kv{$+2.1$} & \kv{$+0.1$} \\\hline
	\end{tabular}
	\end{center}

		\begin{solution}[0.6in]
		Off: $20\log_{10}\sigma$. On: $10\log_{10}(1 + \sigma^2)$, which is
		$10\log_{10} 1.25 = 0.97$ dB and $10\log_{10} 1.64 = 2.15$ dB.
		\end{solution}

	\begin{minipage}[t]{0.58\linewidth}
		\part Averaging $N$ samples divides the floor's spread by $\sqrt N$, so one standard
		deviation above and below the floor is
		\[ 10\log_{10}\left(1 \pm \tfrac{1}{\sqrt N}\right)\ \text{dB}. \]
		Find it for $N = 32$ and $N = 128$, and the gap between the two levels at
		$\sigma = 0.8$.
		\begin{solution}[1.2in]
		\eq{ N = 32:&\quad 10\log_{10}(1 \pm 0.177) = +0.71,\ -0.85\db \\
		N = 128:&\quad 10\log_{10}(1 \pm 0.088) = +0.37,\ -0.40\db \\
		2.15 - (-1.94) &= 4.1\db }
		\end{solution}
	\end{minipage}\hfill%
	\begin{minipage}[t]{0.37\linewidth}
		\raggedleft \vspace{12pt}
		$N=32$: \ansbox[1.0in]{$\approx \pm 0.8$ dB}

		\vspace{4pt}
		$N=128$: \ansbox[1.0in]{$\approx \pm 0.4$ dB}

		\vspace{4pt}
		gap: \ansbox[1.0in]{4.1 dB}
	\end{minipage}

	\part At 15 words per minute, predict the length of a dit, a dah, the gap inside a
	letter, the gap between letters, and the gap after the word, and the number of
	samples in one unit at 32 kHz.

	\begin{center}
	\renewcommand{\arraystretch}{1.25}
	\begin{tabular}{| K{0.62in} | K{0.62in} | K{0.75in} | K{0.75in} | K{0.75in} | K{0.75in} |}
	\hline
	Dit & Dah & Inside a letter & Between letters & After the word & Samples per unit \\\hline
	\kv{80 ms} & \kv{240 ms} & \kv{80 ms} & \kv{240 ms} & \kv{560 ms} & \kv{2560} \\\hline
	\end{tabular}
	\end{center}

		\begin{solution}[0.4in]
		One unit is $1.2/15 = 80$ ms, and $0.080 \times 32{,}000 = 2560$ samples.
		\end{solution}

	\begin{minipage}[t]{0.58\linewidth}
		\part Averaging 32 samples, how many detector outputs is a dit? Averaging 128, how
		long does one output last, and how many outputs is a dit? How many marks (key-downs)
		does one sending of ``ECE'' hold, and how many arrive in 30 s?
		\begin{solution}[1.2in]
		$32/32{,}000 = 1$ ms, so a dit is 80 outputs; $128/32{,}000 = 4$ ms, so a dit is
		20 outputs. ECE is a dit, then dah dit dah dit, then a dit: 6 marks in 26 units,
		2.08 s. In 30 s the flowgraph sends 14 whole sendings, 84 marks, and 4 more marks of
		the 15th, so the count reads about 87 to 88.
		\end{solution}
	\end{minipage}\hfill%
	\begin{minipage}[t]{0.37\linewidth}
		\raggedleft \vspace{12pt}
		$N=32$: \ansbox[1.0in]{80}

		\vspace{4pt}
		$N=128$: \ansbox[1.0in]{4 ms, 20}

		\vspace{4pt}
		marks: \ansbox[1.0in]{6; 87--88}
	\end{minipage}

\end{parts}

%%%%%%%%%%%%%%%%%%%%%%%%%%%%%%%%%%%%%%%%%%%%%%%%%%%
\newpage
\fullwidth{\large\textbf{Part 2: At the Bench} \normalsize(in class)}

\question \textbf{Build and test gr-ece448.} Follow the lab page's steps.

\begin{parts}

	\part Paste the output of \texttt{ctest}, and the \texttt{gr\_python\_install} line
	of \texttt{python/ece448/CMakeLists.txt}.
		\begin{solution}[1.2in]
		\texttt{ctest} reports that 100\% of tests passed, 0 failed out of 1: the one test
		file holds the four test methods of \texttt{qa\_energy\_detector.py}. The line reads\\
		\texttt{gr\_python\_install(FILES \_\_init\_\_.py energy\_detector.py}\\
		\texttt{\ \ DESTINATION ...)}.
		\end{solution}

	\part Open \texttt{L11-keyed-carrier.grc}, which uses the installed block, and run it.
	Paste a screenshot of the Energy Detector's properties dialog.
		\begin{solution}[1.3in]
		The dialog shows Samples Averaged 32, Threshold (dB) \texttt{threshold}, and
		Hysteresis (dB) \texttt{hyst}. On the Embedded Python Block route, the dialog of
		the Embedded Python Block shows the same three parameters, and part (a) comes from
		the lab machine.
		\end{solution}

\end{parts}

%%%%%%%%%%%%%%%%%%%%%%%%%%%%%%%%%%%%%%%%%%%%%%%%%%%
\question \textbf{Set the levels.} In the running flowgraph, set the noise amplitude to
0.5. Then stop it, set \texttt{wpm} to 15 and the two time sinks to 2080 points, one
sending, and run it again.

\begin{parts}

	\begin{minipage}[t]{0.58\linewidth}
		\part Read the floor and the carrier-on level from the power plot, and compare
		them with Question 2(a).
		\begin{solution}[0.8in]
		The floor sits near $-6$ dB with about $\pm 0.8$ dB of spread, and the carrier
		near $+1$ dB, as predicted.
		\end{solution}
	\end{minipage}\hfill%
	\begin{minipage}[t]{0.37\linewidth}
		\raggedleft \vspace{12pt}
		off: \ansbox[1.0in]{$\approx -6$ dB}

		\vspace{4pt}
		on: \ansbox[1.0in]{$\approx +1$ dB}
	\end{minipage}

	\part Set the threshold to the floor's level with no hysteresis, and describe the
	key. Then choose your own threshold and hysteresis so that the key follows the
	keying for at least 20 s. Record your settings and paste a screenshot.
		\begin{solution}[1.3in]
		At $-6$ dB and no hysteresis the key chatters through every gap. A threshold of
		$-2.5$ dB with 1 to 3 dB of hysteresis gives a key that matches the keying: the
		on level sits 3.5 dB above the floor and the off level stays above it.
		\end{solution}

\end{parts}

%%%%%%%%%%%%%%%%%%%%%%%%%%%%%%%%%%%%%%%%%%%%%%%%%%%
\newpage
\question \textbf{Write the Key Timer.} Make a module of your own and write a block in it
from the template, as the lab page lists. The block is a sink with one float input, the
Energy Detector's key, and one parameter, \texttt{rate}, the key's sample rate. It
counts how many samples each mark and each space lasts and keeps the length of the last
mark and the last space, in ms, and the number of marks so far. Three methods,
\texttt{last\_mark()}, \texttt{last\_space()}, and \texttt{mark\_count()}, return them.

\begin{parts}

	\part Paste your \texttt{key\_timer.py}.
		\begin{solution}[2.4in]
		A \texttt{gr.sync\_block} with \texttt{in\_sig=[np.float32]} and
		\texttt{out\_sig=None}. \texttt{work()} loops over the input, compares each value
		(as 0 or 1) with the stored state, and increments a counter while they match. On a
		change, it converts the count to milliseconds,
		\[ t = \frac{1000 \times \text{count}}{\text{rate}}, \]
		stores it
		as the last mark if the state was 1 (and adds one to the mark count) or as the last
		space if it was 0, then stores the new state and restarts the count at 1. The state
		and the count are attributes of \texttt{self}, so a mark that spans two calls to
		\texttt{work()} still gets the right length. \texttt{work()} returns the number of inputs.
		Grade against the faculty reference.
		\end{solution}

	\part Paste your \texttt{qa\_key\_timer.py} test and the output of \texttt{ctest}. The
	test must feed a known key, such as five 0s, three 1s, two 0s, one 1, and four 0s at
	a rate of 1000, and check all three methods.
		\begin{solution}[1.4in]
		For that key the block reports 2 marks, a last mark of 1 ms, and a last space of
		2 ms. The 2 ms space replaces the first space, 5 ms, and the final 4 ms space never
		ends, so the block does not count it. \texttt{ctest} reports the test passed.
		\end{solution}

	\part Paste your \texttt{grc/*.block.yml}, and the \texttt{gr\_python\_install}
	line, which must list \texttt{key\_timer.py}.
		\begin{solution}[1.2in]
		The file has \texttt{make: <module>.key\_timer(\$\{rate\})}, one parameter \texttt{rate} of type
		\texttt{real}, one \texttt{float} input, no outputs. On GNU Radio 3.10.9 the
		install line needs \texttt{key\_timer.py} added by hand.
		\end{solution}

\end{parts}

%%%%%%%%%%%%%%%%%%%%%%%%%%%%%%%%%%%%%%%%%%%%%%%%%%%
\newpage
\question \textbf{Time the Morse.} Save a copy of \texttt{L11-keyed-carrier.grc} with your
settings from Question 4, add your Key Timer on the Energy Detector's key output with
\texttt{rate} = \texttt{samp\_rate / decim}, and show its three values on QT GUI Labels
through Function Probes.

\begin{parts}

	\part Watch the labels for a few sendings and record the mark and space lengths you
	see. Compare them with Question 2(c).

	\begin{center}
	\renewcommand{\arraystretch}{1.25}
	\begin{tabular}{| l | K{0.72in} | K{0.72in} | K{0.72in} | K{0.72in} | K{0.72in} |}
	\hline
	& Dit & Dah & Inside a letter & Between letters & After the word \\\hline
	Measured (ms) & \kv{80--82} & \kv{240--242} & \kv{78--80} & \kv{238--240} & \kv{558--560} \\\hline
	\end{tabular}
	\end{center}

		\begin{solution}[0.5in]
		Each length falls within 1 or 2 ms of 80, 240, 80, 240, and 560 ms.
		\end{solution}

	\part Paste a screenshot of your flowgraph's canvas and of its running window.
		\begin{solution}[1.6in]
		The Key Timer connects to the Energy Detector's key output, and the labels read, for
		example, 80.0 ms and 240.0 ms.
		\end{solution}

\end{parts}

%%%%%%%%%%%%%%%%%%%%%%%%%%%%%%%%%%%%%%%%%%%%%%%%%%%
\newpage
\question \textbf{A noisier signal.} Set the noise amplitude to 0.8.

\begin{parts}

	\part With 32 samples averaged, choose the best threshold and hysteresis you can
	find. Record them, then restart the flowgraph and record the mark count after 30 s.
		\begin{solution}[1.0in]
		The best settings sit near a threshold of $+0.5$ dB with 2 dB of hysteresis, and
		even there the count runs a few marks high, about 89 to 95 for 87 to 88 sent, because
		the noise now and then makes a mark of 1 to 3 ms in a gap. At a threshold of 0 dB the
		count rises to about 130. The 4 dB gap is only about five times the 0.77 dB spread.
		\end{solution}

	\part Stop the flowgraph, set \texttt{decim} to 128, and repeat. Record your settings,
	the count after 30 s, and the dit length the Key Timer reports.
		\begin{solution}[1.0in]
		With a threshold of 0.5 to 1 dB and 1.5 to 2 dB of hysteresis the count matches the
		marks sent, 87 to 88, and the dit reads 80 ms, sometimes 84 ms, one 4 ms output
		long.
		\end{solution}

\end{parts}

%%%%%%%%%%%%%%%%%%%%%%%%%%%%%%%%%%%%%%%%%%%%%%%%%%%
\newpage
\fullwidth{\large\textbf{Part 3: Writeup} \normalsize(individual, after the lab)}

\question \textbf{Hysteresis and timing.} With $\sigma = 0.5$, a $-2.5$ dB threshold, and
3 dB of hysteresis, some dits read 81 or 82 ms; with no hysteresis, they read 80 ms.
Explain why, and what Question 7 says about choosing $N$.
		\begin{solution}[1.5in]
		With 3 dB of hysteresis the off level is $-5.5$ dB, only 0.5 dB above the $-6$ dB
		floor, so after the carrier stops, the averaged noise can stay above the off level
		for an output or two and the key turns off late. A larger $N$ shrinks the floor's
		spread by $\sqrt N$, which lets the levels sit further from the noise, but each
		output then lasts longer, $N/f_s$, so the edges move in larger steps: averaging
		more samples gives a cleaner key with coarser edge timing.
		\end{solution}

\question \textbf{Module or Embedded Python Block?} Lesson 12 builds a Morse decoder that
uses the Energy Detector and your Key Timer. Explain which of the two should live in an
OOT module and why, and what the unit test caught or would catch that watching the
flowgraph would not.
		\begin{solution}[1.5in]
		Both belong in a module: every CW flowgraph in Lesson 12 and later needs them, and a module keeps one
		tested copy of the code. The unit test checks exact values, such as a mark that
		spans two calls to \texttt{work()} or a final space that never ends, which a label
		updating four times a second cannot show.
		\end{solution}

\question \textbf{From timing to letters.} The Key Timer reports marks of 80 and 240 ms at
15 WPM. Propose a rule a decoder can use to tell a dit from a dah, and a rule to tell
the end of a letter from the gap inside one, that still works at 20 WPM.
		\begin{solution}[1.5in]
		Measure the unit from the shortest common mark, then call a mark longer than 2 units
		a dah and a shorter one a dit; call a space longer than 2 units the end of a letter,
		and longer than 5 units the end of a word. Because the rule uses the measured unit,
		60 ms at 20 WPM, it does not depend on the sending speed.
		\end{solution}

\vspace{1.5em}
\noindent\textbf{Documentation:} \rule[-2pt]{0.6\linewidth}{0.4pt}

\end{questions}

\end{document}
